\section{The fourfold and its family}
\label{sec:V-results}

We use the complex orientation and column homology vectors.  A marked
Kuhn cusp means the specified periodic unimodular triangulation, its
integral lattice marking, and its face identifications; these data are
part of the filling.  The two finite fibres have multiplicities five
and four.  The regular fibre is a complex three-torus.

\subsection{The integral eight-manifold}
\RungZero{The displayed rank-six monodromy and the specified affine
fillings produce a three-connected eight-manifold with integral middle
form $3\langle1\rangle\oplus\langle-1\rangle$.  The primitive cusp
collapses five fibre generators, while the actual order-four boundary
map kills each remaining order-two torsion class with coefficient one.}
\RungOne
\begin{enumerate}[label=\arabic*.]
\item The matrices $M_1=A^t$, $M_2=B^t$ have orders five and four, and
  $N=(M_1M_2)^{-1}-I$ has square zero and primitive image of rank three.
  The displayed unimodular marking takes $N$ to
  $\left(\begin{smallmatrix}0&I_3\\0&0\end{smallmatrix}\right)$;
  the invariant covector $\ell(x)=x_0$ proves freeness of every
  nonidentity affine power
  (\cref{maps:src:F.2.5,night:P1.1}).
\item The cusp and finite quotient charts glue to a closed smooth
  manifold.  Killing the cusp angular lattice leaves the central fibre
  generator $h$ and the relations $5a=h$, $4b=-h$, $a+b=0$;
  their integral relation matrix is unimodular, so the fundamental
  group vanishes (\cref{maps:src:F.2.5,night:P1.10}).
\item In the order-four cap the boundary base coordinate is
  $\theta=a+4k$.  Its primitive transverse loop is $(1,0)$ and its
  fibre loop is $(-4,1)$.  The four-prism pullback carries a torsion
  representative $w$ to $(w,-(I+T+T^2+T^3)w)=(w,0)$, so the actual
  transgression is $[w]$, with no extra factor four
  (\cref{mapsaudit:order-four-torsion-transgression}).
\item The three torsion transgressions are isomorphisms onto
  $\mathbb Z/2$.  Together with the primitive free transgression,
  the integral Leray complex gives
  $H_k=\mathbb Z$ for $k=0,8$, $H_4=\mathbb Z^4$, and zero otherwise;
  Hurewicz therefore gives three-connectivity
  (\cref{maps:src:F.2.6,night:P1.12,night:P1.19}).
\item The compatible almost complex structure has $c_1=c_3=0$.
  Todd integrality and the signature bound $|\sigma|\leq4$ force
  $c_2^2=2$, $p_1^2=8$, $p_2=14$ and $\sigma=2$.
  Its unimodular middle lattice is consequently the odd form
  $I_{3,1}$ (\cref{night:P1.13}).
\end{enumerate}

\subsection{The genus-two period construction}
\RungZero{The genus-two pencil
$y^2=x^5+\lambda^3x+\lambda^4$ gives the required rank-three period
planes.  The actual morphism is the projection of
$\ell_0\wedge\mathcal H_{\mathrm{dR}}$, followed by a holomorphic
splitting of the rank-six extension; a scalar shift removes its real
intersection with its conjugate.}
\RungOne
\begin{enumerate}[label=\arabic*.]
\item The pencil has one node at $\lambda_*=-256/3125$ and good
  reductions after root orders eight and five.  Its Hodge bundle on
  that root curve is $T^3\oplus T$, with
  $T=\mathcal O((2/5)[\infty]-(3/8)[0])$ and
  $\deg T=1/40$ (\cref{night:P1.4}).
\item In the exterior-square lattice, projection induces
  $j:\omega^\perp\to(\bigwedge^2\mathbb Z^4)/\mathbb Z\omega$
  with Smith factors $(1,1,1,1,2)$.  This exact index-two map identifies
  the rational systems, and the quotient is the dual of the primitive
  integral lattice.  The lower five-dimensional blocks of the given
  rank-six tuple are the quotient blocks under this map
  (\cref{night:P1.2}).
\item Put $P=p(\ell_0\wedge\mathcal H_{\mathrm{dR}})$, where
  $\ell_0=\langle dx/y\rangle$ and $p$ is projection modulo $\omega$.
  The wedge map has rank three and its image meets $\mathbb C\omega$
  trivially: a decomposable two-vector has square zero, whereas
  $\omega^2\ne0$.  This proves the required rank before a lift is
  chosen (\cref{night:P1.5,maps:src:F.3.5}).
\item The orbifold extension over $P$ splits because
  $H^1(P^*)=0$, as computed from its two line-bundle extensions and
  their integral floor degrees.  A chosen splitting $\sigma$ and
  the canonical functional $\alpha:P\to\mathcal O$ give
  $F_\xi=(\sigma+\xi\alpha e)P$
  (\cref{maps:src:F.3.5,night:P1.7}).
\item On the real line $P\cap\overline P$, choose $n$ with
  $\alpha(n)=1$ and write
  $\sigma(n)-\overline{\sigma(n)}=2ih e$.
  Then $F_\xi\cap\overline{F_\xi}=0$ exactly when
  $h+\operatorname{Im}\xi\ne0$.  At the cusp the period matrix is
  $Z=(\log q)/(2\pi i)I_3+C(q)$, with $C$ holomorphic and bounded;
  hence $h\to+\infty$ and a sufficiently large imaginary shift
  works on the whole completed base
  (\cref{maps:src:F.3.6,night:P1.8}).
\end{enumerate}

\subsection{The complex structure on the connected sum}
\RungZero{The period planes extend over the two finite root charts and
the marked cusp, giving a compact complex fourfold diffeomorphic to
$3\mathbb{HP}^2\#(-\mathbb{HP}^2)$.  Its characteristic vector is
carried explicitly to $(1,1,1,1)$, after which the known smooth
classification and homotopy-sphere absorption identify the oriented
manifold.}
\RungOne
\begin{enumerate}[label=\arabic*.]
\item The lifted periods commute with the finite root actions and
  have the prescribed bounded cusp correction, so the same quotient
  and toroidal charts are holomorphic and proper.  Their collar
  identifications are the marked maps used in the smooth construction
  (\cref{maps:src:F.3.7,night:P1.9}).
\item The preceding integral calculation gives the intersection
  lattice $I_{3,1}$.  The spin characteristic vector $p_1/2$ has
  square two, by the Chern-number calculation
  (\cref{night:P1.13}).
\item In an odd diagonal basis write its coordinates, after sign
  changes, as positive odd integers $(a,b,c,d)$.  A reflection in
  $(1,1,1,1)$ replaces $d$ by $2d-a-b-c$, whose absolute value is
  smaller when $d>1$.  Iteration terminates at $(1,1,1,1)$, proving
  an isometry of the lattice together with its characteristic vector
  (\cref{maps:src:F.4.1}).
\item Wall's oriented almost-diffeomorphism classification, in the
  precise extended-form statement cited there, identifies this pair
  with that of $3\mathbb{HP}^2\#(-\mathbb{HP}^2)$ up to a homotopy
  eight-sphere.  Kramer--Stolz's Theorem A absorbs that sphere in one
  $\mathbb{HP}^2$ summand, giving the stated oriented diffeomorphism
  (\cref{maps:src:F.4.1}).
\end{enumerate}

\subsection{The exact family, its invariants, and its ends}
For this paragraph $\sigma_b$ is the specified algebraic splitting,
and the marking and the Kuhn fan remain those of the construction.
Set
\[
 c_*=-h_{b,\min}=1.0945402503114027801162682630549457037\ldots,
 \qquad \mathcal U=\{\xi\in\mathbb C:\operatorname{Im}\xi>c_*\}.
\]
The exact definition of $c_*$ is the Frobenius period expression in
\cref{night:P1.18}; the decimal is a certified enclosure, not its
definition.  The comparison with the point-normalized parameter is
the full complex translation $\xi_b=\xi_{\rm original}-a_{\rm split}$.

\RungZero{The marked family is parametrized exactly by $\mathcal U$:
its members are pairwise nonisomorphic, have algebraic dimension one
and automorphism group $\mathbb C\times\mathbb Z/2$, and admit no
antiholomorphic isomorphism for this fan.  The real-intersection
function determines the parameter domain; the invariant line and
one half-period determine the automorphisms, while the actual
reflected cusp map fails to preserve a ray of the chosen fan.}
\RungOne
\begin{enumerate}[label=\arabic*.]
\item The ball-arithmetic proof covers the two finite root discs,
  the intervening collar and the cusp annulus, and proves that $h_b$
  has its unique minimum at the infinity root point.  The local
  positive Hessian proves strictness there; the remaining rectangles
  have strict certified lower bounds.  Thus the transversality test
  gives exactly $\mathcal U$ (\cref{night:P1.18}).
\item On the effective root curve of orders five and four put
  $L=\mathcal O((4/5)[\infty]-(3/4)[0])$.  The direct images are
  $R^qh_*\mathcal O=\bigwedge^q
  (\mathcal O\oplus L^{-1}\oplus L^{-2})$, and
  $K_M=h^*L^{-6}$.  The toric normalization complex proves the
  exterior-algebra assertion at the cusp, and the root-chart
  invariant cohomology proves it at the finite fibres
  (\cref{night:P1A.1}).
\item The logarithmic cotangent sequence and its explicit residue
  matrices give the Hodge table below.  The tangent sequence gives
  $h^q(T_M)=(1,1,1,2,0)$; differentiating the actual period lift
  gives a nonzero Kodaira--Spencer class, so the Kuranishi germ is
  smooth of dimension one (\cref{night:P1A.2,night:P1A.3}).
\item An effective horizontal divisor contradicts $H^2(M,\mathbb R)=0$,
  so the meromorphic field is the base field $\mathbb C(\lambda)$.
  Every isomorphism therefore preserves the base and its three
  distinguished values.  The simultaneous linear centralizer is
  $\mathbb Q I$; the order-five affine norm excludes its negative
  integral unit.  The translation sheaf has identity component
  $\mathbb C$ and the displayed order-two half-period, proving the
  automorphism group and equality of parameters for isomorphic
  members (\cref{night:P1A.4}).
\item The two integral anti-intertwiners are $\pm R$.  One is excluded
  by the finite cap norm.  For the other, the actual based cusp
  correction acts on rays by $-P$, and takes $e_2$ to $(0,-1,1)$,
  outside the component's rays $\pm(\{0,1\}^3\setminus\{0\})$.
  The corresponding monomial map contracts a divisor, contradicting
  a biholomorphism (\cref{night:P1.17}).
\item A proposed proper Hausdorff filling across a boundary point of
  a full parameter disc recovers the intrinsic base from $K_M^{-10}$.
  Its marked periods would extend as $F_\xi$ across values where
  $h_b+\operatorname{Im}\xi=0$, contradicting smooth torus fibres
  (\cref{night:P1A.5}).  For a family on a punctured disc whose
  fibres belong to $\mathcal U$, the unique classifying map is
  holomorphic; its Cayley transform extends and the maximum principle
  keeps its limit inside $\mathcal U$.  After the explicit degree-two
  removal of the finite torsor, that family has a smooth filling
  (\cref{night:P1A.6}).
\end{enumerate}

With rows indexed by $p$ and columns by $q$, the Hodge table is
\[
\begin{array}{c|ccccc}
 h^{p,q}&0&1&2&3&4\\\hline
 0&1&1&0&0&0\\
 1&0&2&1&0&0\\
 2&0&1&6&1&0\\
 3&0&0&1&2&0\\
 4&0&0&0&1&1
\end{array}.
\]

\subsection{The effective input classification and the period decision}
\RungZero{Every supplied rank-six integral datum with its affine
orders, twists, clutch and marked primitive cusp has a terminating
conjugacy and integral-topology calculation.  A different, explicit
period-domain construction gives holomorphic transverse periods with
any prescribed finite cusp jet on the punctured base; the polarized
abelian-surface source is excluded by its actual two-line filtration
and the Euler-parity equation.}
\RungOne
\begin{enumerate}[label=\arabic*.]
\item Primitive square-zero cusp matrices have the explicit integral
  normal form and centralizer
  $\left(\begin{smallmatrix}P&Q\\0&P\end{smallmatrix}\right)$.
  Eigenvalue one is necessary for a free affine lift, leaving the
  finite cyclotomic list and its complex-compatibility filter
  (\cref{maps:src:F.2.1,maps:src:F.2.2}).
\item The rational intertwiner algebra, its radical and its
  semisimple quotient reduce the simultaneous integral-conjugacy
  problem to the cited lattice-isomorphism algorithm in its verified
  hypothesis.  Searching by height and lexicographic order terminates
  by the input height, yielding a canonical representative and an
  explicit unimodular conjugator (\cref{maps:src:F.2.36,maps:src:F.2.37}).
\item For that conjugator $R$, the coordinate map $x\mapsto Rx$
  transports $(M_i,v_i,\kappa,P)$ to
  $(RM_iR^{-1},Rv_i,R\kappa,RP)$.  Since
  $(RP)^{-1}(Rx)=P^{-1}x$, it intertwines the specified cusp charts
  as well as the affine caps and clutch.  The resulting finite
  cellular cone and group presentation compute all integral homology
  and the full fundamental group (\cref{maps:src:F.2.31,maps:src:F.2.35,maps:src:F.2.39}).
\item The exact two-point survivor tests and the three-point
  orbifold quotient give the complete constructive parametrization;
  the separate involution branch has 56 classes.  The displayed
  unbounded trace family proves why the terminating bound depends on
  the supplied input (\cref{maps:src:F.2.30,maps:src:F.2.38,maps:src:F.2.39}).
\item Nine explicitly given complex-line directions and a uniform
  singular-value bound put a spanning tube in the unpolarized period
  domain.  The verified local-spray approximation theorem, with the
  relative extension-point argument proved here, produces a
  holomorphic transverse period section and retains any prescribed
  bounded cusp jet (\cref{maps:src:F.3.8,maps:src:F.3.10,maps:src:F.3.11}).
\item For a polarized abelian-surface source, semisimplicity and the
  maximal square-zero cusp force the middle rank-four system to have
  no invariant or coinvariant vectors.  The actual filtration
  $0\subset\mathbb Q\subset E\subset T$, with middle quotient that
  system and top quotient $\mathbb Q$, then forces both relevant
  invariant ranks to be one; the integral Euler-parity relation
  excludes $b_1=b_2=0$ (\cref{maps:src:F.3.4}).
\end{enumerate}
